\section{Liquid-water imbibition: coefficient basis and numerical evidence}
\label{sup:finitecontact}
This section derives the water-transport coefficients of the particle and
gives the numerical evidence for the DT calculation with liquid-water
imbibition. Imbibition is part of the particle formulation and
is active only while the external water film contacts the dry shell. It is
inactive in the Faner experiment, which keeps the reference pore-gas
and retained-water fluxes. Its external series resistance, film-stock limit and energy transport
are given in the main text.
It introduces no extra water storage or retention capacity.

\paragraph{Material interpretation.}
Liquid-water imbibition into soy protein isolates combines adsorption,
capillarity and loosely held interstitial water, and its capacity exceeds
equilibrium vapour sorption several-fold \citep{jovanovich2003wateruptake}.
Assigning the incoming liquid to one retained-water inventory is an
effective closure, not a measurement of its microscopic phase partition. Displacement of hexane by water, swelling, liquid
pressure and blocking of hexane by water are outside this closure.

\paragraph{Reference transport coefficients.}
The reference binary pore value $D_b=6.0\times10^{-10}$ m$^2$/s is an
engineering sensitivity point, 1.5 times the scalar diffusivity used by
\citet{cardarelli2002modeling}. The reference retained-water mobility is derived
from the linearized coupled balances at $T=385.65$ K, $P=130$ kPa and
$y_{\mathrm h}=0.80$, with the reference retention properties.
Eliminating the total molar flux gives the local apparent response
\[
D_{\mathrm{app}}=D_{\mathrm{app,pore}}+g_{\mathrm{ret}}\mathcal L_{\mathrm{w,ret}},
\]
where $D_{\mathrm{app,pore}}=4.54\times10^{-12}$ m$^2$/s and
$g_{\mathrm{ret}}=7.26\times10^{-4}$ m$^3$/mol are obtained from the
same storage and chemical-potential derivatives. Setting the combined
response to $6.2\times10^{-10}$ m$^2$/s gives the reference retained-water
mobility $8.47\times10^{-7}$ mol/(m s). The target value is the maximum apparent
moisture diffusivity reported by \citet{touffet2026moisture} for
pelleted feed at \SI{60}{\celsius}; transferring it to this reference
state is an engineering assumption. On the isothermal, uncapped Luikov
branch, the diffusivity of the retained-water path is
\[
D_{\mathrm{ret}}=
\frac{M_{\mathrm w}\mathcal L A_1}{\rho_{\mathrm{dm}}A_2W^2},
\]
which for the reference mobility is $3.8\times10^{-10}$ m$^2$/s at
$W=0.05$ and $2.4\times10^{-11}$ m$^2$/s at $W=0.20$. That span covers
the two available measurements on oilseed meal: vapour-sorption diffusivities of
sunflower meal, $10^{-12}$ to $1.15\times10^{-10}$ m$^2$/s at
50--\SI{95}{\celsius} \citep{cardarelli1998thesis}, and the soaking
diffusivity of defatted soybean meal, $1.14$ to $2.07\times10^{-11}$
m$^2$/s at 20--\SI{60}{\celsius} \citep{mukherjee2019soaking}. The
reference retained-water mobility is therefore consistent with oilseed meal, although
its value was taken from the pellet study.

\paragraph{The liquid-water mobility factor.}
The liquid-water mobility is $\mathcal L_\ell=999\,\mathcal
L_{\mathrm{w,ret}}=8.47\times10^{-4}$ mol/(m s), so the total shell
mobility is the reference value times a mobility factor of 1000. The
number is the span of the liquid transport
coefficient of porous hygroscopic solids between the dry state and
free water saturation, which \citet{kunzel1995simultaneous} measured on
bricks, sandstones and concretes and expressed as
$1000^{\,w/w_f-1}$ with $w_f$ the free saturation (his eq.~24). In the
sorption region of the same materials the liquid contribution is only
one to two times the vapour transport, so a constant factor is the
water-filled-pore limit of that law. We apply it because the pathway is
active only in the dry shell while the external water film is filling
its pores. Starting from the reference retained-water mobility, $D_{\mathrm{ret}}$ at $W=0.20$ is
$2.4\times10^{-8}$ m$^2$/s on the enhanced pathway.

The mobility factor and the liquid contact conductance act in series, with
the interior
term $d/\mathcal L_\ell$ and the surface term $1/K_\ell$, so a single
uptake time fixes only their sum. Table~\ref{tab:finitecontact:factor}
gives the interior share of the series resistance at the DT radius with
$d=\Rp/5$, the lumped-sphere equivalent, for $K_\ell=0.32$ mol/(m$^2$ s).
Above a factor of about 300 the surface controls the uptake and the
result no longer depends on the factor. Below 100 the interior would
double the resistance, and the conductance derived below would have to
rise to reproduce the same measurement. Completed DT marches at factors
100, 1000 and 10\,000 with the same $K_\ell$
(Table~\ref{tab:finitecontact:factormarch}) confirm this. The tenfold
reduction to 100 lowers the outlet moisture by 0.58 percentage points
and raises hexane by 22.8 mg/kg dry; the tenfold increase to
10\,000 moves them by only $+0.07$ points and $-2.7$ mg/kg dry. The
outlet is on its plateau at 1000.

\begin{table}[htbp]\centering\small
\caption{Interior share of the series resistance to liquid entry at the
DT radius, $K_\ell=0.32$ mol/(m$^2$ s), $d=\Rp/5$.}
\label{tab:finitecontact:factor}
\begin{tabular}{@{}rrr@{}}\toprule
Mobility factor & Interior share & Uptake slowdown vs.\ contact alone\\\midrule
30 & 0.79 & 4.8\\
100 & 0.53 & 2.1\\
300 & 0.27 & 1.4\\
1000 & 0.10 & 1.1\\
3000 & 0.04 & 1.04\\
\bottomrule\end{tabular}\end{table}

\begin{table}[htbp]\centering\small
\caption{DT outlet against the liquid-water mobility factor at
$K_\ell=0.32$ mol/(m$^2$ s); every other input as in
Table~\ref{tab:finitecontact:sensitivity}. The last column gives the
maximum normalized water, hexane and energy defects of the independent
reconstruction from dry-shell formation to exit (acceptance $10^{-10}$).}
\label{tab:finitecontact:factormarch}
\begin{tabular}{@{}rrrrrrl@{}}\toprule
Factor & Moisture, \% wet & Hexane, mg/kg dry & Hexane, ppm wet & Mean $T$, $^\circ$C & Film consumed, s & W / H / U defects\\\midrule
100 & 17.62 & 158.0 & 130.1 & 105.01 & 702.5 & $7.1\times10^{-12}$ / $2.1\times10^{-12}$ / $2.7\times10^{-12}$\\
1000 (adopted) & 18.20 & 135.2 & 110.6 & 104.98 & 624.5 & $4.2\times10^{-11}$ / $8.6\times10^{-13}$ / $6.2\times10^{-12}$\\
10\,000 & 18.27 & 132.5 & 108.3 & 104.97 & 614.0 & $2.6\times10^{-11}$ / $2.8\times10^{-12}$ / $4.4\times10^{-12}$\\
\bottomrule\end{tabular}\end{table}

\paragraph{Construction of the contact conductance.}
\label{sup:contactscale}
The liquid contact conductance represents the resistance to liquid entry
between the external water film and the connected particle matrix. It keeps
the surface conductance finite when the shell forms, where the main-text law
gives $G_\ell\to K_\ell$ as $d\to0$. Its value is derived from the only kinetic
measurement of solvent-extracted soybean meal in liquid water
\citep{hemmingsen2008water}: the coarse sieve fraction (above 0.5 mm,
mean 1.43 mm; 25 g fat per kg dry matter; initial water 11.8\%, that is
$W_0=0.134$ kg/kg dry) immersed at about \SI{80}{\celsius}: its
rehydration index is on its plateau at the first sample, 90 s, and
does not rise afterwards. The construction assumes a sphere of radius
$R=0.715$ mm, entry through the surface alone, the surface at the film's state
($W_{\mathrm{cap}}$) and the retention law of the main text. The water
balance $(\rho_{\mathrm{dm}}V/M_{\mathrm w})\dot W
=-A K_\ell[\psi(W)-\psi(W_{\mathrm{cap}})]$ with
$\psi(W)=(1-A_1/W)/A_2$ integrates in closed form to
\[
t=\frac{\rho_{\mathrm{dm}}R A_2}{3M_{\mathrm w}A_1K_\ell}\,
W_{\mathrm{cap}}\!\left[(W_0-W_1)+W_{\mathrm{cap}}
\ln\frac{W_{\mathrm{cap}}-W_0}{W_{\mathrm{cap}}-W_1}\right],
\]
and requiring $W_1=W_{\mathrm{cap}}-0.01$ kg/kg at $t=90$ s gives
$K_\ell=0.253$ mol/(m$^2$ s) at \SI{80}{\celsius}. Scaling the entry
rate with the inverse viscosity of water from 80 to \SI{100}{\celsius}
(a factor of 1.26) gives 0.318, adopted as 0.32. The end-point tolerance
is the largest single uncertainty of the construction: 0.02 or 0.005
kg/kg below capacity gives 0.21 or 0.43. Reading the measurement at
\SI{20}{\celsius} with the same construction and viscosity scaling
(a factor of 3.55) gives 0.90, so the adopted value is the lower,
less extrapolated, of the two readings.

The construction assumes that the sieve size is a sphere diameter, that
the interior had reached the retention capacity when the first sample was
drained, and that the entry rate scales with the liquid viscosity. The
measurement is on a drained filter cake and does not separate particle,
interstitial and surface water. The 90 s is therefore an upper bound on the
entry time, and the derived conductance is the slowest value consistent
with the measurement. It is a literature-derived value under these
declared assumptions, not a measurement on DT meal, and no DT outlet value
entered the calculation.

With the adopted liquid-water mobility, the resistance length
$\mathcal L_\ell/K_\ell$ is 2.65 mm, 1.76 times the DT radius. At
$d=\Rp$ the series conductance is $G_\ell=0.204$ mol/(m$^2$ s), and the
contact term's share of the series resistance is 0.64. The uniform-sphere linearization
$\tau_K=\rho_{\mathrm{dm}}\Rp A_2W^2/(3M_{\mathrm w}K_\ell A_1)$ gives
56 s at $W=0.20$, and the closed-form fill of a DT sphere limited by contact alone,
from the feed loading 0.124 to $W_{\mathrm{cap}}-0.01$, takes 153 s
(491 s in the earlier reference case 0.10). These are scale estimates;
radial gradients, pore-gas exchange, the changing surface state, the
film's water stock and the capacity all remain active in the calculation.

\paragraph{Sensitivity to the contact conductance.}
\label{sup:contactsensitivity}
Table~\ref{tab:finitecontact:sensitivity} lists completed DT marches at
four liquid contact conductances with every other input held constant,
the earlier reference case 0.10 among them. The outlet moisture is bounded by the condensate supply
and the retention capacity: the external water film is consumed before
the exit in every case, and a fourfold change in $K_\ell$ moves the outlet
moisture by 1.3 percentage points and the wet-basis hexane by 43 ppm.
The conductance sets how long the film persists, not how much water the
particle finally holds.

\begin{table}[htbp]\centering\small
\caption{DT outlet against the liquid contact conductance; mobility
factor 1000, two radial cells, regime-C maximum step 5 s, exit at
1766.5 s. Wet-basis values exclude the external film, which is zero at
the exit in every case.}
\label{tab:finitecontact:sensitivity}
\begin{tabular}{@{}rrrrrr@{}}\toprule
$K_\ell$, mol/(m$^2$ s) & Moisture, \% wet & Hexane, mg/kg dry & Hexane, ppm wet & Mean $T$, $^\circ$C & Film consumed, s\\\midrule
0.05 & 16.68 & 194.5 & 162.0 & 105.09 & 816.8\\
0.10 & 17.40 & 165.3 & 136.5 & 105.03 & 733.6\\
0.20 & 17.94 & 144.8 & 118.8 & 104.99 & 662.4\\
0.32 (adopted) & 18.20 & 135.2 & 110.6 & 104.98 & 624.5\\
\bottomrule\end{tabular}\end{table}

\paragraph{Events and refinement.}
The adopted DT calculation uses the feed history up to 492.257073 s,
effective radius 1.5 mm, two radial cells, the prescribed bath and the
existing retention capacity. The moving-front calculation begins with a
0.01-s shell-formation interval, advances regime B in 0.1-s intervals,
shortened where events require, and resolves grid-face arrival and departure
before core extinction at 503.560810 s. External water stays positive
throughout regime B, so the zero-film regime-B branch, which is not
implemented, is not exercised (Sec.~\ref{sup:dtmarch}).
Regime C advances the stationary particle and the film together to the
film's depletion at 624.548745 s and then to the exit.
Table~\ref{tab:finitecontact:refinement} gives the regime-C step
refinement, performed in the earlier reference case $K_\ell=0.10$ from the
same moving-front end point; the 5-s maximum step is used in every
reported march. The extra digits resolve the numerical sensitivity; the main
text uses rounded values.

\begin{table}[htbp]\centering\small
\caption{Regime-C timestep sensitivity of the DT particle with liquid-water
imbibition, $K_\ell=0.10$ mol/(m$^2$ s).}
\label{tab:finitecontact:refinement}
\begin{tabular}{@{}lrr@{}}\toprule
Quantity & Maximum step 10 s & Maximum step 5 s\\\midrule
Moisture, \% wet basis & 17.3813 & 17.4033\\
Hexane, ppm wet basis & 137.222 & 136.474\\
Mean temperature, $^\circ$C & 105.0329 & 105.0314\\
External-water depletion time, s & 735.395 & 733.575\\
\bottomrule\end{tabular}\end{table}

\paragraph{Independent balances.}
An independent reconstruction of the adopted march, from dry-shell
formation to exit, evaluates water (W), hexane (H) and energy (U) in the
particle and external water film at every accepted step. The maximum
normalized W/H/U defects are $4.24\times10^{-11}$, $8.61\times10^{-13}$ and $6.16\times10^{-12}$,
below the $10^{-10}$ acceptance limit; the regime-B segment alone gives
$1.05\times10^{-11}$, $8.61\times10^{-13}$ and $3.94\times10^{-12}$. The
unchanged feed history keeps its earlier check. Component partial
enthalpies count each energy transfer once. These checks show conservation in the
declared numerical experiment;
whole-history mesh and time-step convergence remains open
(Sec.~\ref{sup:dtmarch}), and conservation does not identify the material
coefficients.
